% Generated draft: requires mathematical and source review.
\documentclass[11pt]{article}
\usepackage[margin=25mm]{geometry}
\usepackage{fontspec}
\setmainfont{Latin Modern Roman}
\usepackage{amsmath,amssymb,mathtools,unicode-math}
\setmathfont{Latin Modern Math}
\usepackage{enumitem,needspace,xurl,hyperref}
\setlength{\parindent}{0pt}
\setlength{\parskip}{5pt}
\setlength{\emergencystretch}{3em}
\newcommand{\N}{\mathbb{N}}
\newcommand{\Z}{\mathbb{Z}}
\newcommand{\Q}{\mathbb{Q}}
\newcommand{\R}{\mathbb{R}}
\newcommand{\C}{\mathbb{C}}
\newcommand{\F}{\mathbb{F}}
\newcommand{\A}{\mathbb{A}}
\newcommand{\PP}{\mathbb{P}}
\newcommand{\E}{\mathbb{E}}
\newcommand{\eps}{\varepsilon}
\newcommand{\dd}{\,\mathrm{d}}
\begin{document}
\section*{UnsolvedMath ID: 30004066}
\textbf{Source problem code:} catalogue entry
\subsection{Definitions and mathematical statement}
Binary tomography concerns reconstruction of an object whose values belong to a prescribed set of two gray levels. The supplied record does not specify these levels, the reconstruction domain, the measurement operator, the data, the objective function, or the constraints of the referenced primal problem. It also does not specify the Lagrangian or the admissible dual variables.

Let \(p^\star\) and \(d^\star\) denote, schematically, the optimal values of the binary-tomography primal problem and its Lagrangian dual referenced by the source. Strong duality asks whether
\[
 p^\star=d^\star.
\]
This equality expresses equality of optimal values; it does not by itself assert attainment. The source record does not provide the hypotheses or quantifiers under which this equality is being asked, nor its conventions for infeasible problems or infinite optimal values. Consequently, this is a schematic statement of the target, not a fully specified optimization problem.

The original source record poses three questions:
\begin{enumerate}
\item Does strong duality hold for the referenced binary-tomography primal and dual formulations?
\item Can the preceding unique-solution/intersection conjecture be proved? The conjecture itself is absent from the supplied record: neither the objects being intersected nor the conditions relating their intersection to uniqueness are stated. Its mathematical content therefore cannot be reconstructed precisely.
\item Can the approach be extended to discrete tomography with a prescribed finite set \(\Gamma\) of gray levels satisfying \(\lvert\Gamma\rvert>2\), by developing an analogous dual formulation? The intended primal formulation and the properties required of an analogous dual formulation are unspecified.
\end{enumerate}
The cleaned statement retains the first and third questions but omits the conjecture-proof request present in the original statement. All three targets are retained here, with their unresolved qualifications made explicit.
\subsection{Short English statement}
The source asks whether the referenced binary-tomography problem has zero duality gap, whether an earlier uniqueness/intersection conjecture can be proved, and whether the method extends to more than two gray levels. The optimization formulations and the conjecture are missing, so a complete mathematical statement requires further source review.
\subsection{Sources}
\begin{itemize}
\item Supplied local UnsolvedMath registry record, permanent ID \(30004066\), title: Strong Duality and Multilevel Discrete Tomography; exact source problem code: \texttt{catalogue entry}. The record belongs to the Oberwolfach Reports Open Problems collection and gives proposed year 2019; no individual proposer is identified.
\item Supplied report citation: \emph{Tomographic Inverse Problems: Theory and Applications} (2020), Oberwolfach Reports, pp.\ 209--303, DOI \texttt{10.4171/owr/2019/4}. \url{https://doi.org/10.4171/owr/2019/4}. The supplied background describes these reports as records of mathematical workshops and their open-problem discussions. No precise page or excerpt containing the formulations or conjecture is supplied.
\item Supplied literature reference: Kadu--van Leeuwen, convex binary tomography. \url{https://doi.org/10.1109/TCI.2019.2898333}.
\item Literature assessment dated exactly \texttt{2026-08-22}: the supplied curation labels the questions open and states that the source records them as open questions with only small-scale computational evidence. This is supplied dated curation, not an independently verified literature assessment. The registry also labels the statement \texttt{corrected\_verified} and reports checking it against the official Oberwolfach report; that verification claim is supplied metadata and has not been independently verified here.
\end{itemize}
\end{document}
